% ======================================================================
%  Lethe -- research: private AI access, the competitive landscape.
%  Source: reports/Private AI access competitive landscape.md
%  Build from this directory:  pdflatex landscape.tex  (run twice)
% ======================================================================
\documentclass[10pt,a4paper]{article}
\usepackage[a4paper,top=20mm,bottom=24mm,left=20mm,right=20mm,footskip=11mm]{geometry}
\usepackage[T1]{fontenc}
\usepackage[utf8]{inputenc}
\usepackage[type1,mainfont,semibold,proportional]{inter}
\usepackage[type1,scale=0.88]{plex-mono}
\usepackage{microtype}
\usepackage[table]{xcolor}
\usepackage{tikz,amsmath,amssymb}
\usepackage{booktabs,tabularx,array,adjustbox}
\usepackage{flafter}
\usepackage{enumitem}
\usepackage[nobottomtitles*]{titlesec}
\usepackage[font={sf,small},labelfont={bf},labelsep=period,skip=5pt]{caption}
\usepackage[skins,breakable]{tcolorbox}
\usepackage{fancyhdr}
\usepackage[numbers,sort&compress,square]{natbib}
\usepackage{xurl}
\usepackage[hidelinks]{hyperref}
% ======================================================================
%  lethe-iso.tex -- isometric TikZ macro library + palette for Lethe
% ======================================================================
\usetikzlibrary{calc,shapes.geometric,arrows.meta,decorations.pathmorphing,
  decorations.markings,backgrounds,positioning,fit,patterns,patterns.meta,
  shadows.blur,matrix}

% ---------------------------------------------------------------- palette
\definecolor{ink}{HTML}{14161F}
\definecolor{paper}{HTML}{FBFAF6}
\definecolor{indigo}{HTML}{4F46E5}
\definecolor{violet}{HTML}{7C3AED}
\definecolor{teal}{HTML}{0D9488}
\definecolor{river}{HTML}{0EA5E9}
\definecolor{coral}{HTML}{F43F5E}
\definecolor{amber}{HTML}{F59E0B}
\definecolor{lime}{HTML}{65A30D}
\definecolor{slate}{HTML}{64748B}
\definecolor{sky}{HTML}{38BDF8}
\definecolor{plum}{HTML}{A21CAF}
\definecolor{sand}{HTML}{EAE6DA}

% ---------------------------------------------------------------- projection
\tikzset{
  iso/.style={x={(0.866cm,-0.5cm)}, y={(0.866cm,0.5cm)}, z={(0cm,1cm)}},
  isolabel/.style={font=\sffamily\bfseries\tiny, text=ink, align=center, inner sep=1pt},
  isocap/.style={font=\sffamily\tiny, text=ink!70, align=center, inner sep=1pt},
  flow/.style={-{Stealth[length=4pt,width=3.5pt]}, line width=0.9pt, line cap=round},
  flowd/.style={-{Stealth[length=4pt,width=3.5pt]}, line width=0.9pt, dashed, line cap=round},
  faceline/.style={line width=0.3pt, line join=round},
  zone/.style={font=\sffamily\bfseries\tiny, inner sep=1pt},
}

% \isobox{x}{y}{z}{w}{d}{h}{color}{label}
\newcommand{\isobox}[8]{%
  \fill[#7!16!white, draw=#7!60!black, faceline]
    (#1,#2,{#3+#6}) -- ({#1+#4},#2,{#3+#6}) -- ({#1+#4},{#2+#5},{#3+#6}) -- (#1,{#2+#5},{#3+#6}) -- cycle;
  \fill[#7!52!white, draw=#7!60!black, faceline]
    (#1,#2,#3) -- ({#1+#4},#2,#3) -- ({#1+#4},#2,{#3+#6}) -- (#1,#2,{#3+#6}) -- cycle;
  \fill[#7!82!white, draw=#7!60!black, faceline]
    ({#1+#4},#2,#3) -- ({#1+#4},{#2+#5},#3) -- ({#1+#4},{#2+#5},{#3+#6}) -- ({#1+#4},#2,{#3+#6}) -- cycle;
  \node[isolabel] at ({#1+#4/2},{#2+#5/2},{#3+#6}) {#8};
}

% \isoboxc{x}{y}{z}{w}{d}{h}{color}{label}  -- label centred on the front-left face
\newcommand{\isoboxc}[8]{%
  \isobox{#1}{#2}{#3}{#4}{#5}{#6}{#7}{}%
  \node[isolabel, text=white] at ({#1+#4/2},#2,{#3+#6/2}) {#8};
}

% \isoplate{x}{y}{z}{w}{d}{color}  -- ground plane with a faint grid
\newcommand{\isoplate}[6]{%
  \fill[#6!5!white, draw=#6!35, faceline]
    (#1,#2,#3) -- ({#1+#4},#2,#3) -- ({#1+#4},{#2+#5},#3) -- (#1,{#2+#5},#3) -- cycle;
  \begin{scope}
    \clip (#1,#2,#3) -- ({#1+#4},#2,#3) -- ({#1+#4},{#2+#5},#3) -- (#1,{#2+#5},#3) -- cycle;
    \foreach \i in {1,...,40} {
      \draw[#6!11, line width=0.2pt] ({#1+\i},#2,#3) -- ({#1+\i},{#2+#5},#3);
      \draw[#6!11, line width=0.2pt] (#1,{#2+\i},#3) -- ({#1+#4},{#2+\i},#3);
    }
  \end{scope}
}

% \isostack{x}{y}{z}{w}{d}{n}{color}  -- n thin slabs (each 0.22 tall, gap 0.14)
\newcommand{\isostack}[7]{%
  \foreach \k in {0,...,\numexpr#6-1\relax} {
    \isobox{#1}{#2}{#3+\k*0.36}{#4}{#5}{0.22}{#7}{}
  }
}

% \isocube{x}{y}{z}{s}{color}
\newcommand{\isocube}[5]{\isobox{#1}{#2}{#3}{#4}{#4}{#4}{#5}{}}

% \isocyl{x}{y}{z}{r}{h}{color}{label}  -- vault cylinder (drawn in 2D at projected centre)
\newcommand{\isocyl}[7]{%
  \path (#1,#2,#3) coordinate (cylb);
  \path (#1,#2,{#3+#5}) coordinate (cylt);
  \fill[#6!70!white, draw=#6!60!black, faceline]
    ($(cylb)+(-#4,0)$) arc (180:360:{#4} and {#4*0.5}) -- ($(cylt)+(#4,0)$) arc (0:-180:{#4} and {#4*0.5}) -- cycle;
  \fill[#6!16!white, draw=#6!60!black, faceline] (cylt) ellipse ({#4} and {#4*0.5});
  \node[isolabel] at (cylt) {#7};
}

% \isoriver{x}{y1}{y2}{w}  -- a band along y between x and x+w on the z=0 plane
\newcommand{\isoriver}[4]{%
  \fill[river!35!white, draw=none]
    (#1,#2,0) -- ({#1+#4},#2,0) -- ({#1+#4},#3,0) -- (#1,#3,0) -- cycle;
  \foreach \o in {0.2,0.5,0.8} {
    \draw[river!80, line width=0.5pt, decorate, decoration={snake, amplitude=0.6pt, segment length=6pt}]
      ({#1+#4*\o},#2,0) -- ({#1+#4*\o},#3,0);
  }
}

% small legend chip:  \chip{color}{text}
\newcommand{\chip}[2]{\tikz[baseline=-0.5ex]{\fill[#1,rounded corners=1pt] (0,-0.09) rectangle (0.24,0.09);}\,{\sffamily\scriptsize #2}}


\renewcommand{\familydefault}{\sfdefault}
\linespread{1.18}\raggedbottom\frenchspacing
\widowpenalty=10000 \clubpenalty=10000 \displaywidowpenalty=10000
\setlength{\parindent}{0pt}\setlength{\parskip}{5.5pt plus 1.5pt minus 0.5pt}
\setlist{nosep,leftmargin=1.2em,itemsep=2pt,topsep=3pt}
\newcommand{\eyebrow}[1]{{\fontsize{7}{9}\selectfont\bfseries\textcolor{indigo}{\MakeUppercase{#1}}}}
\urlstyle{tt}\renewcommand{\UrlFont}{\ttfamily\color{ink!62}}
\setcitestyle{square,numbers,comma}

% ---------------------------------------------------------------- headings
\renewcommand{\thesection}{\arabic{section}}
\newcommand{\two}{\ifnum\value{section}<10 0\fi}
\titleformat{\section}
  {\sffamily\bfseries\fontsize{13}{16}\selectfont\color{ink}\raggedright\nohyph}
  {\textcolor{indigo}{\two\thesection}}{0.6em}{}
\titlespacing*{\section}{0pt}{18pt plus 4pt minus 2pt}{7pt}
\titleformat{\subsection}
  {\sffamily\bfseries\fontsize{10.5}{13}\selectfont\color{ink}\raggedright\nohyph}
  {\textcolor{indigo}{\two\thesection.\arabic{subsection}}}{0.6em}{}
\titlespacing*{\subsection}{0pt}{12pt plus 3pt minus 1pt}{4pt}
\renewcommand{\bottomtitlespace}{0.12\textheight}

% ---------------------------------------------------------------- callouts
\newtcolorbox{callout}[2][indigo]{
  enhanced jigsaw, breakable, colback=#1!4!white, colframe=#1!4!white, boxrule=0pt, arc=1.5pt,
  left=9pt, right=9pt, top=6pt, bottom=7pt, before skip=8pt, after skip=10pt,
  borderline west={2.5pt}{0pt}{#1},
  fontupper=\small\linespread{1.16}\selectfont,
  title={{\fontsize{7}{9}\selectfont\bfseries\textcolor{#1!85!black}{\MakeUppercase{#2}}}},
  attach title to upper={\\[3pt]}}

% ---------------------------------------------------------------- tables
\captionsetup[table]{position=top,skip=4pt}
\newcommand{\tablefont}{\footnotesize\linespread{1.08}\selectfont\setlength{\tabcolsep}{4pt}%
  \renewcommand{\arraystretch}{1.22}\rowcolors{2}{white}{slate!6}}
\newcolumntype{L}[1]{>{\raggedright\arraybackslash\hsize=#1\hsize}X}
\newcolumntype{B}[1]{>{\raggedright\arraybackslash\bfseries\hsize=#1\hsize}X}
\newcommand{\colhead}[1]{\textcolor{ink!70}{\bfseries #1}}
\newcommand{\tablenote}[1]{\par\vspace{4pt}{\scriptsize\color{ink!70}#1\par}}
\setlength{\textfloatsep}{14pt plus 3pt minus 2pt}
\setlength{\floatsep}{12pt plus 3pt minus 2pt}
\renewcommand{\topfraction}{0.9}\renewcommand{\bottomfraction}{0.6}
\renewcommand{\textfraction}{0.08}\renewcommand{\floatpagefraction}{0.8}
\setcounter{topnumber}{2}\setcounter{bottomnumber}{1}\setcounter{totalnumber}{3}

% ---------------------------------------------------------------- page style
\pagestyle{fancy}\fancyhf{}
\renewcommand{\headrulewidth}{0pt}
\fancyfoot[L]{\fontsize{7}{9}\selectfont\textcolor{ink!55}{LETHE\,·\,research\,·\,October 2026}}
\fancyfoot[R]{\fontsize{7}{9}\selectfont\textcolor{ink!55}{\thepage}}

% ---------------------------------------------------------------- references
\renewcommand{\bibsection}{\section*{References}\phantomsection\addcontentsline{toc}{section}{References}}
\renewcommand{\bibfont}{\footnotesize\linespread{1.05}\selectfont}
\setlength{\bibsep}{1.6pt plus 0.5pt}

\hypersetup{pdftitle={Private AI access: the competitive landscape},
  pdfauthor={Pranjal Joshi, Prashil Bulbule}}

\begin{document}

% ======================================================================  title
\thispagestyle{fancy}
\vspace*{-3mm}
\eyebrow{Lethe · research · October 2026}\\[7pt]
{\fontsize{24}{28}\selectfont\bfseries\color{ink}Private AI access:\\ the competitive landscape}\\[4pt]
{\fontsize{14}{17}\selectfont\color{ink!80}Consumers buy anonymous AI; funds buy custody}\\[7pt]
{\fontsize{9}{11}\selectfont\color{ink!65}Pranjal Joshi\quad·\quad Prashil Bulbule}
\vspace{3mm}

\begin{callout}{Summary}
As of October 2026, almost every layer of private AI access is crowded. The one empty spot is the exact combination Lethe proposes, and for institutions it is empty because no institution has asked for it. More than a dozen products already pool many users behind one provider account in front of GPT, Claude and Gemini, including Duck.ai, Brave Leo, Kagi, Venice, 1Privacy, nullsink and Apple's ChatGPT extension. Stanford's OpenAnonymity already puts Privacy Pass blind signatures in front of those models, for free. Every enclave startup (Tinfoil, Confer, Privatemode, Confident Security, Maple) and every crypto network (NEAR AI, Phala, Nillion, Secret) can prove what happens to the prompt, but all of them serve only open-weight models. Apple, Google and Meta already run anonymous tokens, oblivious relays and attested compute for hundreds of millions of their own users. Nobody sells one product that gives closed frontier models four things together: unlinkability from the operator enforced by cryptography, oblivious transport, a published anonymity-set size, and a contract the model provider has signed. Consumer demand for privacy-framed frontier AI is proven: Venice is profitable at more than \$70M ARR and was valued at \$1B. Institutional demand for hiding a firm's identity from the model lab is not. JPMorgan, Goldman Sachs, Morgan Stanley, Balyasny and Freshfields already run OpenAI and Anthropic models through channels where the lab never sees their prompts. When Anthropic imposed 30-day retention in 2026, the CISOs of the big banks negotiated logs held in their own clouds, not anonymity. As pitched, with \$100k--500k firm memberships as the profit centre and a break-even consumer app as ``the crowd'', Lethe would not be funded by a seed investor or by YC. The profit centre has no observed buyer, and it runs into both the labs' new rules on detecting misuse across sessions and regulators' rules on traceability. The consumer half gives away the only part with proven demand. And the claimed moat (``a competitor needs a bigger crowd'') belongs to incumbents who already have millions of users. The company worth funding turns the plan around. It sells a consumer and prosumer product priced like AI and run for profit. It is built on ``accountable anonymity'' that a lab can sign a Duck.ai-style contract for. Institutional access is reframed as provider-blind, firm-audited custody and kept as a later option.
\end{callout}

\begin{figure}[!b]
\centering
\figfit{\textwidth}{figs/positioning}
\caption{\textbf{Frontier relays rest on policy promises; only research systems reach the cryptographic corner.} Horizontal: which models a product serves. Vertical: what its operator can still learn about its users. For frontier-model requests only OpenAnonymity, zkAPI and LLMTor make users cryptographically unlinkable, and none of them combines that with commercial quality and a contract the provider has signed. Positions are qualitative, read from the text of this report.}
\label{fig:map}
\end{figure}
\section{More than a dozen products already pool users in front of frontier labs}

The anonymous-access market falls into four groups. A UMass Amherst study lumps most of them together as proxies that ``exchange content confidentiality for operational practicality''~\cite{arxiv-2609-33903-umass}. \textbf{Big-brand relays} (Duck.ai, Brave Leo, Kagi Assistant, Apple's signed-out ChatGPT extension) strip IP and account data and send requests from the company's own provider accounts under contract. \textbf{Crypto-paid, no-account API proxies} (nullsink, PayPerQ, nadanada, NanoGPT, AnonRouter, 1Privacy) do the same with prepaid balances. \textbf{Open-weight and enclave apps} (Proton Lumo, Maple, Confer, Venice's Private and TEE tiers) avoid the frontier labs entirely. \textbf{Research systems} (OpenAnonymity, zkAPI, LLMTor) add blind-signature or zero-knowledge unlinkability. \textbf{Pooling, then, is not a differentiator: every product in the first two groups already makes the lab see one customer instead of many.} The line that matters is whether the operator can tell its own users apart. For frontier-model requests, only OpenAnonymity, zkAPI and LLMTor make that cryptographically impossible. Brave's Premium tokens unlink payment from usage, but the scheme is not documented.

\textbf{OpenAnonymity is the closest existing implementation of Lethe's consumer layer.} It was built by Ken Liu (Stanford AI Lab, advised by Percy Liang and Sanmi Koyejo) and Erik Chi (University of Michigan, advised by J. Alex Halderman). Users get blind-signed ``inference tickets'' made with Cloudflare's open-source privacypass-ts library. A ``station'' checks a ticket without learning who holds it and issues a short-lived child API key on a shared OpenRouter account. The browser then calls OpenRouter directly~\cite{openanonymity-oa-station-r,openanonymity-oa-chat-priv}. A verifier running in an AMD SEV-SNP enclave audits each station's OpenRouter privacy settings and looks for shadow accounts with logging turned on~\cite{openanonymity-oa-verifier}. Its claim is narrow and honest. Prompts are not hidden from the provider, only made unlinkable, so an observer ``cannot tell if 100 requests came from 100 people with 1 request each, or 1 person with 100 requests''~\cite{openanonymity-oa-chat-gith}. Its limits are stated just as plainly. The browser proxy hides the user's IP from the station, but not necessarily from OpenRouter. The command-line client exposes the source IP by default and reuses one credential for 60 seconds. Billing ``intentionally exposes subscriber IP''. And IP, fingerprinting and DNS ``remain exposed without additional privacy infrastructure''~\cite{openanonymity-oa-chat-priv}.

The project has served \textbf{more than 15,000 prompts at Stanford}~\cite{ken-liu-on-x}. It is ``currently university-funded'' and is weighing foundation grants, venture money or a nonprofit structure~\cite{brave-technologist-podcast}. On 1 October 2026 the Ethereum Foundation and the project launched \textbf{zkAPI}. Users deposit ETH or USDC and then prove each request is paid for without revealing which deposit paid~\cite{the-block-zkapi-launch}. Its roadmap is a ``secretary model'': a local agent holds personal data and releases only what each remote model needs. A memory repository, nanomem, is already public~\cite{brave-technologist-podcast,openanonymity-on-github}. That is very nearly Lethe's consumer memory app. The project's own prompt scrubber shows the ceiling: in the UMass study, \textbf{50--75\% of in-set SWE-Chat users} were still re-identified in at least one conversation despite text-rewriting defences~\cite{arxiv-2609-33903-umass}.

\begin{table}[tbp]
\centering\tablefont
\caption{Pooled and anonymous access products in front of frontier and open-weight models.}\label{tab:access}
\begin{tabularx}{\textwidth}{@{}B{1.064}L{1.392}L{0.942}L{1.105}L{0.696}L{0.901}L{0.901}@{}}
\toprule
\colhead{Product} & \colhead{Mechanism and trust model} & \colhead{Frontier models} & \colhead{Identity, payment, network} & \colhead{Price} & \colhead{Funding, traction} & \colhead{Heading} \\
\midrule
OpenAnonymity & Privacy Pass tickets exchanged for short-lived child keys on a shared OpenRouter account; SEV-SNP verifier & GPT, Claude, Gemini via OpenRouter & Operator cannot link users (cryptographic); IP partly exposed; zkAPI ETH/USDC & Not public & University-funded; $>$15k prompts & Local ``secretary'' agent, memory, zkAPI \\
Venice & Proxy (Anonymous mode), own GPUs, NEAR and Phala enclaves & GPT, Claude, Gemini, Grok in Anonymous mode & Account, IP and timestamps processed; card or crypto & \$0--200/mo & \$65M at \$1B; $>$\$70M ARR; $\sim$3M users & Own data centres \\
Duck.ai & IP-stripping proxy; zero retention by contract & GPT-5.x, Claude Opus 4.8 & No account; card for paid tiers & Free; \$9.99; \$19.99 & DuckDuckGo; no AI-specific numbers & Paid tiers \\
Brave Leo & Reverse proxy; unlinkable Premium tokens; NEAR enclave pilot & Claude Haiku/\allowbreak Sonnet/\allowbreak Opus & Token scheme undocumented; IP hiding is policy & Free; \$15/mo & Brave & Client-side attestation \\
Kagi Assistant & Account-based; no account data sent to providers & GPT 6.x, Claude 5.5 and Fable 5.1, Gemini 3.x & Privacy Pass covers Search only & Premium models on Ultimate & Kagi & Privacy Pass for Assistant \\
Proton Lumo & Zero-access history; open-weight models on Proton servers & None & Free without account & Free; \$12.99 & Proton & Memory, images, search \\
1Privacy & Relay strips identifiers and rotates keys & 350+ models & Recovery email only, no identity checks & From \$25/mo & Undisclosed & ``AI Identity'', ``AI Cleanup'' \\
nullsink & No-log prepaid gateway & Anthropic, OpenAI, Tinfoil & XMR, BTC, Lightning & Model rates plus fee & Operator undisclosed & Not stated \\
\bottomrule
\end{tabularx}
\end{table}

\textbf{Venice is the commercial proof that consumers will pay, and its ``Anonymous'' mode is already a pool.} In July 2026 it raised a \textbf{\$65M Series A at a \$1B valuation}, its first outside capital, led by Dragonfly with Coinbase Ventures, F-Prime, Archetype and others. At the time it had \textbf{more than \$70M of annualised revenue, profitable since Q1 2026}, about 3M users and 1.7M API calls a day~\cite{the-block-venice-funding,techcrunch-venice-series-a}. One third-party estimate puts ARR closer to \$60M~\cite{ownyourmind-venice}. CEO Erik Voorhees (ShapeShift) and president and CTO Jesse Proudman run about 45 staff, and they are spending the round on their own GPUs and data centres~\cite{siliconangle-venice}. Venice offers four modes~\cite{venice-privacy}. Anonymous mode proxies GPT, Claude and Gemini with the user's identity hidden, but tells users to ``assume the provider is storing your content''. Private mode runs on Venice's GPUs or zero-retention partners. TEE and E2EE modes run on NEAR AI and Phala enclaves and require Pro. The trust model is policy, not cryptography: Venice processes ``account identifiers, API key identifiers, timestamps, model selection, token counts, and IP addresses''~\cite{venice-docs-privacy}. Pricing is free for 10 prompts a day, then \$18 for Pro, \$68 for Pro Plus and \$200 for Max per month~\cite{venice-pricing}. Its VVV and DIEM tokens turn staked capital into API credit. Two confounders weaken the lesson: Venice sells itself as ``private and uncensored'', and about 8\% of users pay in crypto~\cite{techcrunch-venice-series-a}.

The big brands own distribution. \textbf{Duck.ai} strips IP and metadata so requests ``appear as though they are coming from DuckDuckGo''. It has agreements with OpenAI, Anthropic, Mistral, Azure and Tinfoil covering zero retention, no training and no sharing, and its help page lists GPT-5.6 Sol and Terra and Claude Opus 4.8, Sonnet 4.6 and Haiku 4.5~\cite{duckduckgo-help-duck-ai}. An older version of its terms gives providers up to 30 days to delete~\cite{duck-ai-privacy-terms}. It is free, with Plus at \$9.99 and Pro at \$19.99 a month~\cite{spread-privacy-duckduckgo}. \textbf{Brave Leo} routes requests through a reverse proxy, discards responses, and gives Premium subscribers ``unlinkable tokens so that payment details cannot be associated with usage''. It hosts Claude Haiku, Sonnet and Opus alongside Qwen and Llama, for \$15 a month~\cite{wikipedia-brave-leo}. It has piloted NEAR AI enclaves for DeepSeek V3.1 in its Nightly build, with attestation checked on Brave's servers rather than on the user's device~\cite{brave-blog-browser-ai}. \textbf{Kagi Assistant} shares no account information with providers and offers GPT 6.x, Claude Opus 5.5 and Fable 5.1, and Gemini 3.x~\cite{kagi-docs-assistant}. Its Privacy Pass tokens, shipped in February 2025, still cover \textbf{Search only}; Assistant is ``soon to follow''~\cite{kagi-docs-privacy-pass}. \textbf{Proton Lumo 2.0} runs only open-weight Qwen 3.5 and GLM 5.2 on Proton servers, with zero-access chat history. It can be used free without an account, and Plus costs \$12.99~\cite{macrumors-proton-lumo-2}. \textbf{Apple's ChatGPT extension}, used signed out, hides the user's IP and bars OpenAI from storing or training on requests, though it still sends time zone, country, device type and language~\cite{apple-chatgpt-extension-pr}.

The long tail sells the same promise more crudely. \textbf{1Privacy} (1Defender Inc.) strips account, IP and session data at a relay and rotates API keys ``frequently enough that traffic-pattern analysis can't cluster requests''. It offers 350+ models from \$25 a month, with no identity checks and no audit~\cite{1privacy-ai}. \textbf{nullsink} takes Monero, Bitcoin and Lightning and stores only a hash of the API key and the balance. It reaches 54 models from Anthropic, OpenAI and Tinfoil using its own credentials, and it admits that it ``reads request content in process'' and that ``you trust the nullsink operator''~\cite{nullsink-trust}. \textbf{AnonRouter} offers 215+ models in Anonymous, Private, Intel TDX and end-to-end-encrypted tiers, for a 5--8\% fee~\cite{anonrouter-ai}. \textbf{PayPerQ}, founded in February 2024 by Matt Ahlborg, needs no registration and averages about 2 cents a query~\cite{ppq-ai-payperq}. \textbf{LLMTor} redeems blind RSA tokens over Tor, and its builder concedes that banning abusers needs ``one permanent blind RSA token per user'', which brings linkability back~\cite{hacker-news-llmtor}.

Taken together, these products leave the same holes. None of the relays uses Oblivious HTTP, none publishes how many users share each upstream account, and none pads or batches traffic. Every frontier relay can see each user's prompts in plaintext and is trusted not to log them. Abuse control under anonymity is either undocumented or, as at LLMTor, solved by giving up anonymity. \textbf{Lethe's statement that nobody has assembled the primitives in front of a frontier model is therefore false for the individual pieces.} Blind signatures in front of GPT, Claude and Gemini ship today at OpenAnonymity and LLMTor, unlinkable payment tokens ship at Brave, and zero-knowledge payments ship at zkAPI. The statement holds only for the full stack, under contract, at commercial quality.

\section{Enclave startups prove what happens to the prompt, but only on open-weight models}

Confidential-inference startups protect something different: what the prompt says. They protect it better than any relay can, but only for models whose weights they hold.

\textbf{Tinfoil (YC X25)} runs open-weight models on NVIDIA confidential GPUs behind AMD SEV-SNP. It measures the firmware, kernel, initrd, configuration and model weights (via dm-verity), publishes the measurements to a Sigstore transparency log, and has its client SDKs check on every connection that the server's TLS key is the attested one~\cite{tinfoil-docs-attestation-a}. The founders are Tanya Verma (from Cloudflare's cryptography team), Jules Drean (MIT, trusted hardware), Sacha (MIT, privacy-preserving cryptography) and Nate (Tor). They launched on HN in May 2025 at \$20 a month and \$2 per million API tokens~\cite{launch-hn-tinfoil}. In September 2026 Tinfoil wrote that ``all our hosted models are open-weight and widely available via other providers''. It also moved its safety classifiers inside the enclave so that ``the category that was violated'' stays private~\cite{tinfoil-blog-safety-withou}. Its own privacy policy shows the metadata gap. It collects names, emails and sign-in data via Clerk, billing data via Stripe and the app stores, IP addresses, approximate region and per-request usage~\cite{tinfoil-privacy-policy}. Funding is unclear: \$125K from YC, Pioneer Fund, and an undisclosed May 2026 round~\cite{pitchbook-tinfoil}. No revenue or customer counts are public. Its direction is a bridge to closed weights. Tinfoil Containers (March 2026) host ``proprietary models'' in attested VMs on up to eight H200 or B200 GPUs~\cite{tinfoil-blog-tinfoil-conta}. A June 2026 ``verifiable clean room'' ran evaluations on GLM-5.1 without the weights leaving the enclave~\cite{tinfoil-blog-verifiable-cl}. No lab is named as a partner.

\textbf{Confer}, Moxie Marlinspike's AI assistant, has the strictest consumer design found. Encryption keys are derived from the user's passkey. Inference runs in a confidential VM whose entire root filesystem is measured with dm-verity. Builds are reproducible with Nix and mkosi and published to a Sigstore log, and the attestation is bound into the Noise Pipes handshake between client and server~\cite{confer-blog-private-infere}. It is free for 20 messages a day and \$35 a month for unlimited use, on open-weight models it does not name~\cite{techcrunch-confer}. In March 2026 Marlinspike announced that Confer's technology will ``underpin Meta AI'' while Confer stays independent~\cite{confer-blog-meta-ai}. In September it added Confidential Workers, attested Linux VMs that accept a single connection, for tool use~\cite{confer-blog-confidential-w}. No anonymous payment or IP hiding was found, and funding is undisclosed.

\textbf{Privatemode} (by Edgeless Systems, formerly Continuum) is the German, EU-hosted option for government and regulated buyers. It runs open-weight GLM-5.3 and gpt-oss-120b in H100 and B200 confidential environments at Scaleway and Lyceum, and a client-side proxy verifies the workers before any keys are exchanged~\cite{privatemode-docs-architect}. At \textbf{\texteuro{}0.65 per million output tokens} for gpt-oss-120b it has the cheapest published price~\cite{privatemode-pricing}. Edgeless has disclosed no funding round since a \$5M seed in 2023~\cite{tech-eu-edgeless-systems}. \textbf{Confident Security} is the one startup built to hide users at the network level. Its CONFSEC product copies Apple's Private Cloud Compute: Oblivious HTTP through a third-party relay, GPU attestation and transparency logs. It was open-sourced in November 2025 as OpenPCC, built by former Databricks and Apple engineers under CEO Jonathan Mortensen~\cite{business-wire-openpcc-laun,finsmes-confident-security}. It raised a \$4.2M seed, which a later release puts at \$5M, and it serves Mistral, Gemma, Llama, Qwen and DeepSeek, not GPT or Claude~\cite{techcrunch-confident-secur,confident-security}. \textbf{Maple} (built on OpenSecret) goes furthest on identity. It offers signup without an email, which co-founder Mark Suman calls ``quite popular'', Bitcoin payment, and an OpenAI-compatible API running in enclaves. Models include gpt-oss-120b, Kimi K2.6, GLM-5.1 and DeepSeek V4 Pro, at \$0.90--3.00 per million input tokens~\cite{privacy-guides-forum-maple,maple-blog-maple-proxy}. Its backing comes from the Human Rights Foundation rather than a venture round~\cite{human-rights-foundation}. Enterprise money sits one layer up. \textbf{Opaque} raised \$24M at a \$300M valuation in February 2026~\cite{siliconangle-opaque}, naming ServiceNow, Anthropic and Accenture as customers or partners~\cite{pr-newswire-opaque-series}.

\begin{table}[tbp]
\centering\tablefont
\caption{Confidential-inference startups.}\label{tab:enclave}
\begin{tabularx}{\textwidth}{@{}B{0.917}L{1.125}L{0.833}L{1.083}L{0.833}L{1.125}L{1.083}@{}}
\toprule
\colhead{Company} & \colhead{What is attested} & \colhead{Models} & \colhead{Identity, payment, network} & \colhead{Price} & \colhead{Funding, team} & \colhead{Heading} \\
\midrule
Tinfoil (YC X25) & Boot chain and weights; Sigstore log & Open-weight only (stated) & Email/OAuth, Stripe, IP logged & \$20/mo; API pay-as-you-go & YC \$125K plus undisclosed; ex-Cloudflare, MIT, Tor & Containers for proprietary models; memory, search \\
Confer & Whole root filesystem; attestation in the handshake & Open-weight & Passkey; no anonymous payment found & Free; \$35/mo & Undisclosed; Marlinspike & Meta AI integration; agent workers \\
Privatemode & Attested workers (Contrast) & GLM-5.3, gpt-oss & Account and API key & From \texteuro{}0.65/M output & \texteuro{}1.45M plus \$5M (2023) & EU government, regulated sectors \\
Confident Security & OHTTP relay plus GPU attestation & Open-weight & IPs hidden from compute; business billing & Usage-based & \$4.2M--\$5M seed & OpenPCC standard; sells to providers \\
Maple & Enclave inference & Open-weight & No-email accounts; Bitcoin & Free; \$20; \$100 & HRF-backed & Projects; DeepSeek V4 Flash \\
\bottomrule
\end{tabularx}
\end{table}

Two weaknesses run through the whole category. \textbf{Cheap physical attacks broke the hardware guarantees in late 2025.} Battering RAM used a device costing about \$50 to break SGX and SEV-SNP~\cite{the-hacker-news-battering}. TEE.fail used under \$1,000 of hardware against DDR5 memory to extract TDX and SEV-SNP attestation keys and to ``borrow'' NVIDIA GPU attestations, and Intel and AMD plan no firmware fix~\cite{tee-fail}. \textbf{Performance is also worst exactly where frontier models live.} Tinfoil's own June 2026 benchmarks show \textbf{118\% added latency between tokens} for vLLM on multi-GPU Blackwell at batch size one. Overhead falls to about 1\% only for single-GPU models on runtimes built for confidential computing~\cite{tinfoil-blog-confidential}. For Lethe the lesson cuts both ways. No enclave startup can serve GPT, Claude or Gemini, so a frontier pool does not compete with them on models. But every one of them positions itself against redaction backed only by promises, and none attests that it does not log identity, billing or IP. Those records sit outside the enclave everywhere.

\section{Crypto networks run the same enclaves, at valuations far ahead of usage}

In 2026, ``decentralised private inference'' means, almost without exception, open-weight models running in Intel TDX confidential VMs on NVIDIA H100-to-B300 GPUs, with attestation reports and a token attached. \textbf{NEAR AI Cloud} launched in December 2025, from the network co-founded by Illia Polosukhin. It encrypts prompts on the device and serves DeepSeek V4, GLM 5.3, Kimi K2.6/K3 and Qwen 3.x in TDX enclaves, and Brave, OpenMind and Phala were launch customers~\cite{near-ai-blog-near}. In 2026 it added Intel Trust Authority attestation, staking-for-credits, private USDC payments (``Confidential Intents''), a GitHub Copilot extension, and distribution through OpenRouter and SayGm~\cite{near-ai-blog,cryptobriefing-near-ai-git}. It publishes no usage figures, and its OpenRouter route is explicitly not confidential~\cite{near-ai-blog-openrouter}. The 500,000 NEAR staked in the first two weeks yields only about \textbf{\$5,000 a month} of credit under NEAR's own 100:1 rule~\cite{cryptobriefing-near-ai-sta}.

\textbf{Phala}, whose CEO is Marvin Tong according to a weak source~\cite{coindesk-pha}, routes traffic through an attested gateway and then an attested model VM, and it signs a receipt for every response. Confidential models cost \$0.07 to \$3.00 per million input tokens. Phala claims a live counter of about \textbf{33.7 billion confidential tokens a day} and 5,000+ users~\cite{phala,phala-private-ai-inference}. Its RedPill gateway lists GPT, Claude and Gemini as ``protected by a TEE proxy''; the proxy is attested, but the lab is not~\cite{phala-cloud-redpill}. \textbf{Nillion} markets multi-party computation and homomorphic encryption, but its own docs say nilAI runs LLM inference in NVIDIA confidential-computing enclaves, on Gemma 4 and Qwen 3.6~\cite{nillion-docs-llms}. It reports \textbf{1.4M inferences in total} against \$52.7M raised~\cite{nillion-phase-2-upgrade,binance-research-nillion}. \textbf{Secret Network} registers attested GPU workers on-chain and splits user payments according to tokens served~\cite{secret-network-docs-secret}. It is the one network that tries to hide payment metadata, but its SCRT token is down \textbf{96\% in a year} to a market cap of about \$5.9M~\cite{coingecko-api-scrt}. \textbf{Chutes}, a Bittensor subnet, moved fully onto TDX enclaves with optional end-to-end encryption~\cite{chutes}, and Targon runs TDX endpoints on NVIDIA Hopper and Blackwell GPUs~\cite{iq-wiki-targon}. Atoma, Ritual, Fortytwo and Gensyn show no private-inference traction in 2026.

\begin{table}[tbp]
\centering\tablefont
\caption{Decentralised private-inference networks.}\label{tab:crypto}
\begin{tabularx}{\textwidth}{@{}B{0.875}L{1.167}L{0.833}L{1.000}L{1.042}L{1.083}@{}}
\toprule
\colhead{Network} & \colhead{What serves traffic} & \colhead{Closed frontier models} & \colhead{Reported usage} & \colhead{Market cap (7 Oct 2026) or funding} & \colhead{Heading} \\
\midrule
NEAR AI & Single-operator TDX and NVIDIA confidential computing & None in an enclave & Not disclosed & NEAR $\sim$\$6.6B (the main chain, not AI) & Staking credits, agents, IDE and gateway distribution \\
Phala / RedPill & Attested gateway plus TDX GPU VMs & Proxied through an attested gateway & $\sim$33.7B tokens/day (self-reported) & PHA $\sim$\$54M & Ethereum L2, GPU rental, OpenRouter route \\
Nillion & NVIDIA confidential-computing enclave (not MPC) & None & 1.4M inferences in total & NIL $\sim$\$47M; \$52.7M raised & Ethereum L2, Blacklight verification \\
Secret & Enclave workers registered on-chain & None & None found & SCRT $\sim$\$5.9M (down 96\% in a year) & SecretVM on Azure, GCP, AWS \\
Chutes (Bittensor SN64) & TDX and NVIDIA confidential computing & None & $\sim$120B tokens/day, mostly not confidential (undated source) & Subnet token $\sim$\$126M & Full enclave fleet \\
\bottomrule
\end{tabularx}
\tablenote{Market data from CoinGecko PHA~\cite{coingecko-api-pha}, CoinGecko NIL~\cite{coingecko-api-nil} and TAO Media~\cite{tao-media-chutes}.}
\end{table}

Two findings matter for Lethe. First, enclaves do nothing for identity, payment or network metadata. NEAR's staking leaves a permanent on-chain link between a wallet and an AI account, Nillion takes card payments through Stripe, and no network documents Tor, mixnet or OHTTP support. Second, the closest existing analogue to ``a privacy pool in front of GPT and Claude'' is not a token network. It is \textbf{SayGm}, an AI gateway that runs inside a TDX enclave, takes prepaid credits, and has a ``Frontier tier'' that sends Claude and GPT requests upstream, where ``that provider's policy governs your data''~\cite{near-ai-blog-saygm}. SayGm sells nearly Lethe's consumer bundle, minus the pooled unlinkability. It also shows that an attested relay, which makes ``we don't log'' checkable at the intermediary, is now an off-the-shelf technique that Lethe should adopt rather than a moat it can claim.

\section{Apple, Google and Meta already ship Lethe's stack, but only for their own users}

The pattern Lethe proposes is already standard at three of the largest companies in the world: anonymous tokens for sign-in and rate limiting, a third-party oblivious relay, attested compute, and a transparency log. In \textbf{Apple's Private Cloud Compute}, the device gets a Token Granting Token, issued per user and device with RSA blind signatures. It exchanges that token, behind an OHTTP relay, for batches of one-time tokens, and requests go to nodes chosen at random~\cite{wisec-2026-arxiv-2605}. Cloudflare lists PCC as an OHTTP customer~\cite{cloudflare-blog-ohttp-gate}. The WiSec researchers reverse-engineered the system and found three problems. The one-time-token nonce, SHA256(TGT$\Vert$salt), lets requests be linked to each other. The backend skips signature validation and accepts tokens that have been swapped between devices or reused. And Apple's response changed only its documentation~\cite{wisec-2026-arxiv-2605}. Under a January 2026 deal, Gemini-based Apple Foundation Models run inside PCC~\cite{9to5mac-gemini-and-pcc}. Since June 2026 PCC has been moving onto Google Cloud, using NVIDIA confidential GPUs, Intel TDX and two independent roots of trust~\cite{apple-security-research-ex}.

\textbf{Google's Private AI Compute} runs Gemini on TPU enclaves behind IP-blinding relays ``operated by third-parties''. It swaps device credentials for an ``anonymous, cacheable token'' so that a user's identity ``never travels with their data on the inference path''. In September 2026 it announced per-user memory stored on the server inside an enclave, which is persistent and tied to the user by design~\cite{google-private-ai-compute,help-net-security-private}. \textbf{Meta's Private Processing} signs users in with Meta's open-source Anonymous Credential Service. It routes requests through a Fastly OHTTP relay, which also distributes key configurations so that Meta cannot hand individual users their own keys and single them out. Compute runs on SEV-SNP confidential VMs with H100 GPUs, and Cloudflare runs the artifact log. Incognito Chat on WhatsApp and the Meta AI app began rolling out on this system in May 2026~\cite{meta-private-processing-te,meta-newsroom-incognito-ch}. It is an exception mode: by default, Meta has used AI chats to target ads since December 2025~\cite{the-register-meta-ai}.

\begin{table}[tbp]
\centering\tablefont
\caption{Platform and lab positions on anonymous access.}\label{tab:bigtech}
\begin{tabularx}{\textwidth}{@{}B{0.972}L{1.268}L{0.887}L{0.803}L{0.887}L{1.183}@{}}
\toprule
\colhead{Player} & \colhead{Anonymous access today} & \colhead{Whose models} & \colhead{Status} & \colhead{Direction} & \colhead{Risk to a pool} \\
\midrule
Apple PCC and ChatGPT extension & Blind-signature tokens and OHTTP; signed-out ChatGPT, no storage & Apple's own and Gemini-based; ChatGPT & Shipped; linkability flaw unfixed & PCC on Google Cloud & Highest: free and includes a rival's model, but Apple devices only \\
Google Private AI Compute & Anonymous tokens, third-party relays, TPU enclaves & Gemini & Narrow Pixel features & Per-user sealed memory; ``open-minded'' on ads & Medium (consumer) \\
Meta Private Processing and Incognito Chat & Anonymous credentials, Fastly OHTTP relay, confidential VMs & Meta AI & Rolling out & Confer integration; ads remain the default & Medium (consumer) \\
OpenAI & None; Private Safety Processing, Private Inference & GPT & Preview & Confidential but identified & Low for anonymity \\
Anthropic & None; 30-day retention, Enterprise Frontier Safeguards & Claude & In force / phasing in & Monitoring held by the customer & Low; controls access to the top models \\
Azure, Bedrock & Lab can't see data; customer is identified & Many & Generally available & Not researched & Already gives institutions a lab-blind route \\
\bottomrule
\end{tabularx}
\end{table}

The frontier labs are moving the other way, toward \textbf{confidential but identified}. \textbf{OpenAI} began showing ads in February 2026 that may draw on memory and recent chats when memory is on~\cite{openai-help-ads-in}. In August it previewed Private Safety Processing, under which zero-retention customers keep content under their own keys while OpenAI still detects misuse across sessions~\cite{tnw-private-safety-process}. At DevDay on 29 September it announced ``Private Inference'' on confidential computing for enterprises~\cite{venturebeat-openai-private}. \textbf{Anthropic} has, since 9 June 2026, retained prompts and outputs for \textbf{30 days} on its ``Covered Models'' (Fable 5 and 5.1, Mythos 5 and 5.1). This applies even to zero-retention customers, and on Bedrock, Google Cloud and Azure, so that Anthropic can catch misuse ``that may only become apparent across multiple interactions''~\cite{claude-help-covered-models}. Mythos 5 is ``available only to approved partners under limited access programs''~\cite{claude-support-covered-mod}. Anthropic's September answer, Enterprise Frontier Safeguards, stores the monitoring data in the customer's own cloud bucket under the customer's keys, with no Anthropic human review, and Anthropic says it is ``working to support third-party offerings''~\cite{anthropic-enterprise-front}. \textbf{Azure and Bedrock} already keep the lab from seeing customer data. On Bedrock, each provider's model sits in an escrow account with no access to prompts or logs~\cite{aws-re-post-bedrock}. On Azure Foundry, prompts ``are NOT available to OpenAI''. But every request carries the customer's resource ID~\cite{microsoft-learn-azure-open}.

The labs' terms set the boundaries for any pool. OpenAI's Services Agreement lets a customer make ``Customer Applications available to End Users'' (\S{}2.2), but it bars reselling or leasing account access (\S{}3.1), requires each End User Account to belong to one person (\S{}3.2), and forbids transferring API keys (\S{}3.3). It also lets OpenAI require that a specific end user be suspended (\S{}8.1), and it makes the customer warrant export-control compliance for its end users, ``including any `know your customer' requirements'' (\S{}16.11)~\cite{openai-services-agreement}. Anthropic permits powering ``products and services Customer makes available to its own customers and end users'' (\S{}A.1) and bars resale ``except as expressly approved'' (\S{}D.4)~\cite{anthropic-commercial-terms}. Its rules for unsupported regions cover any entity more than 50\% owned from those regions, reportedly including ``cloud service intermediaries''~\cite{anthropic-unsupported-regi}. So a consumer app is permitted. A pool of anonymous firms looks closer to resale, which needs approval. Either way, the pool must hold a per-member handle it can suspend, and it must know enough about institutional members to vouch for their sanctions and ownership status. Google's Gemini and Vertex terms were not researched.

The plumbing underneath is becoming a commodity. Cloudflare launched an OHTTP Gateway in closed beta on 2 October 2026, naming Flo Health's ``Anonymous Mode'' as a customer~\cite{cloudflare-blog-ohttp-gate}. Fastly sells an OHTTP relay~\cite{fastly-docs-oblivious-http}. And authors from Apple and Cloudflare are standardising ARC, an anonymous credential that can be shown a fixed number of times and is designed for rate-limiting anonymous clients~\cite{ietf-draft-privacy-pass}. On this evidence, any lab could add OHTTP and anonymous tokens to its API within a year at modest cost. What stops OpenAI and Anthropic is not engineering but their own safety rule that misuse of frontier models must be detectable across sessions. Platform risk is therefore highest from Apple, which already gives free, anonymous, no-storage access to a rival's model, and lowest from the labs themselves.

\section{YC funded every piece except the pool, and capital followed revenue}

Two independent sweeps of the YC directory found \textbf{no YC company running a pooled-anonymity proxy in front of GPT, Claude or Gemini}. The first, in September, covered 2,398 companies from W23 to F26; the second, in October, covered 1,950 companies from W24 to F26~\cite{yc-oss-yc-companies}. Instead, YC has funded the pieces: enclave infrastructure, content redaction and data-loss prevention, user-owned context, and a crowded lane of agent identity and authorisation.

\begin{table}[tbp]
\centering\tablefont
\caption{YC companies nearest to Lethe.}\label{tab:yc}
\begin{tabularx}{\textwidth}{@{}B{0.958}L{0.375}L{1.583}L{1.250}L{0.833}@{}}
\toprule
\colhead{Company} & \colhead{Batch} & \colhead{What it does} & \colhead{Relation to Lethe} & \colhead{Funding, status} \\
\midrule
Tinfoil & X25 & Verifiable enclave inference, open-weight models only & Same promise, opposite trade-off; calls redaction a band-aid & YC plus Pioneer Fund; active \\
LogosGuard & F25 & Local browser extension that redacts PII before prompts reach ChatGPT, Claude, Gemini; YC profile now reads as AI governance & Hides content, not identity & Two founders, ex-Apple AI; active \\
Egoist Machines & S26 & ``AI Passport'': portable personal context, scoped and revocable access, local hardware planned & Makes identity portable; Lethe makes it unlinkable & Two founders (Cambridge, ex-OpenAI) \\
Unbound & S24 & Enterprise AI gateway against data leakage, since repositioned toward agent access & Gateway that runs under the enterprise's own identity & \$4M seed (Race Capital) \\
Caution & S26 & Attested, reproducible enclave hosting; has ``private AI inference'' customers & Infrastructure Lethe could buy & Team from Distrust, Turnkey, BitGo \\
Credal.ai & W23 & Redaction before cloud LLMs, now ``control plane for enterprise agents'' & Left the privacy wedge & \$4.8M seed (Spark Capital) \\
Blyss & W23 & ``End-to-end encrypted AI'' in GPU enclaves & Tinfoil's thesis, two years early & Inactive \\
Multifactor, Alter, Agentic Fabriq, Clawvisor, Decawork & S25--S26 & Agent identity, scoped credentials, audit trails & The agent-gateway lane, now crowded & Seed \\
Candor, Golf, compliant-llm, Galini & W23--X25 & Shadow-AI data-loss prevention, MCP policy & Privacy sold from the CISO's budget & Seed \\
LiteLLM, Respan, Experiential Labs, Pangolin, Consus & W23--S26 & AI gateways and routers; Pangolin adds identity-aware access, Consus routes ITAR/CUI traffic & The delivery layer, commoditised & LiteLLM \$1.6M seed \\
Sentient OS, RunAnywhere, Cactus, Epicenter, Skillsync & S25--F26 & On-device inference, local-first memory & The ``why not just run it locally?'' alternative & Seed \\
\bottomrule
\end{tabularx}
\tablenote{Sources: Tinfoil~\cite{y-combinator-tinfoil}, LogosGuard~\cite{y-combinator-logosguard,logosguard-site}, Egoist~\cite{y-combinator-egoist-machin}, Unbound~\cite{pulse-2-0-unbound}, Caution~\cite{y-combinator-caution}, Credal~\cite{techcrunch-credal}, Blyss~\cite{y-combinator-blyss}, LiteLLM~\cite{y-combinator-litellm}.}
\end{table}

\begin{figure}[tbp]
\centering
\figfit{0.94\textwidth}{figs/funding}
\caption{\textbf{Capital followed revenue and enterprise buyers.} Disclosed rounds and valuations; figures are as reported in the sources cited in this section. Bar length is the round size; colour is the shape of company. Privacy sold to developers as a cryptographic property raised \$4--5M seeds, while OpenRouter, the neutral gateway in front of every model, shows the ceiling.}
\label{fig:funding}
\end{figure}

The venture record shows three shapes that raised real money and one that did not. Consumer private AI that had revenue before it raised, funded by crypto investors, produced Venice's \$1B round. Enterprise confidential-AI platforms with blue-chip customers produced Opaque's \$24M Series B. Enterprise AI data-leak gateways sold to CISOs produced the largest outcomes. WitnessAI raised a \$58M Series B on ARR growth above 500\%~\cite{axios-witnessai}, SentinelOne bought Prompt Security for about \$133.6M in cash plus stock~\cite{sec-10-q-sentinelone}, and Cato bought Aim Security~\cite{cyberscoop-cato-aim-securi}. The fourth shape, privacy sold to developers as a cryptographic property, stalled at seed or died. Confident Security raised \$4.2M, Unbound \$4M, and Credal \$4.8M before it pivoted. Edgeless has disclosed no round since 2023~\cite{cb-insights-edgeless-syste}. Mithril pivoted out of confidential inference into web agents before it was sold~\cite{techleap-mithril-security}, and Blyss went inactive after YC.

Two neighbouring cases show the ceiling. OpenRouter, the neutral gateway in front of every model, went from raising \$40M at about \$500M in June 2025, to \$113M at \$1.3B in May 2026, to a reported Stripe acquisition above \$7B~\cite{contrary-research-openrout,bloomberg-stripe-and-openr}. And the vertical apps that hold model contracts on behalf of firms earn on capability, not privacy: Harvey has \$400M+ ARR at a \$15.5bn valuation~\cite{legal-business-harvey}, and Rogo has about \$53M ARR across 300+ institutions~\cite{sacra-rogo}.

YC's own signals are mixed but readable. Neither the Spring nor the Fall 2026 Requests for Startups includes privacy-first AI access~\cite{let-s-data-science,explainx-yc-requests-for}, yet the S26 batch still funded Egoist and Caution, so the category is not disqualifying. The objections a partner would raise are already visible in the comments on Tinfoil's Launch HN~\cite{launch-hn-tinfoil}. One commenter wrote, ``Nobody wants you as a service provider. Everyone wants Gemini, OpenAI''. Others said enclaves only move the trust to Intel, AMD and NVIDIA, and asked where the SOC 2 report was. YC's portfolio adds further objections: Venice and DuckDuckGo already do this, the product depends on the labs' terms, consumers do not pay for privacy, and YC watched Blyss die on nearby ground.

\section{Institutions buy custody and audit trails, not anonymity}

Lethe's institutional thesis rests on two premises: that finance and law firms avoid frontier models to keep their names off requests, and that they would pay to hide who is asking.

\textbf{The first premise is out of date.} About 250,000 JPMorgan staff use the bank's LLM Suite, which draws on OpenAI and Anthropic models~\cite{cefpro-jpmorgan-llm-suite}. Goldman has Anthropic engineers embedded to build Claude agents~\cite{hargreaves-lansdown-goldma}. Over 98\% of Morgan Stanley's advisor teams use a GPT-4 assistant under zero data retention~\cite{openai-morgan-stanley}. About 95\% of Balyasny's investment teams use an in-house platform that reasons with GPT-5.4~\cite{cfotech-balyasny}. Freshfields gives Claude to 5,700 staff under a multi-year Anthropic deal~\cite{freshfields-anthropic-part}. Harvey reaches about 80\% of the Am Law 100~\cite{legal-business-harvey}. A Lowenstein Sandler survey of 107 private-equity, hedge-fund and venture firms found 100\% using AI and 96\% raising AI budgets in 2026~\cite{lowenstein-sandler-alterna}. Large firms reach frontier models through internal gateways (JPMorgan's LLM Suite, D.E. Shaw's LLM Gateway, Man Group's bridge), hyperscaler cloud accounts, direct lab partnerships, or vertical vendors that hold the contracts~\cite{resonanz-capital}. On the Bedrock and Azure routes, the lab already cannot see the firm's prompts.

\textbf{For the second premise there is no evidence at all.} Every survey, deployment write-up and regulatory document found frames the risk as content: material non-public information, client data, training, human review and legal discovery. AIMA's survey of hedge-fund managers found the top generative-AI concern was confidential information leaking into external repositories~\cite{aima-hedge-funds-and}. Lowenstein's top concerns were data ownership and privacy (42\%), third-party cyber risk (41\%) and the cost of vetting vendors (40\%), and none of its answer options asked about the provider knowing the firm's interests. Mayer Brown's M\&A guidance treats prompts as discoverable evidence and recommends protecting them by contract~\cite{mayer-brown-m-a}. The cleanest test came in 2026. When Anthropic imposed 30-day retention on its best models, the CISOs of Goldman, Morgan Stanley, Citi, Bank of America and Wells Fargo, working together through the ARC systemic-risk group, co-designed Enterprise Frontier Safeguards. Monitoring data now sits in their own cloud buckets under their own keys, and Anthropic still knows exactly who they are~\cite{anthropic-enterprise-front}. \textbf{They asked for custody, not anonymity.}

The idea that knowing who looked at what is valuable is well supported. Hedge funds identified by their IP addresses in SEC EDGAR logs earn about 1.5\% higher annualised abnormal returns when they pull public filings~\cite{management-science}. Investor-relations analytics de-anonymise institutional visitors to company websites, and one vendor says it has identified 150+ activists, 10\% of them before their campaigns became public~\cite{irwin}. In 2013 Bloomberg reporters were found to be reading client login data~\cite{cnbc-jpmorgan-and-bloomber}. But in the Bloomberg case, clients demanded a policy firewall and stayed; they did not move to anonymous terminals. The dark-pool analogy also turns around on inspection. Off-exchange trading passed 50\% of US equity volume in 2025, with dark pools at about 17.7\%~\cite{cboe-2025-us-equities}. But dark pools are regulated, identified intermediaries that hide a client from counterparties who could trade against it, and the labs are not yet counterparties. The nearest thing to a future motive is the labs moving into finance. OpenAI hired 100+ former bankers to train its models on deal work~\cite{sherwood-news-openai-and}, Anthropic sells Claude for Financial Services directly, and in September 2026 Robinhood put OpenAI and Anthropic agents inside retail order flow~\cite{fortune-robinhood-agents}.

Regulation pushes the same way. FINRA's 2026 oversight report expects generative-AI communications to be retained, prompt and output logs to be stored, and third parties' use of generative AI to be assessed~\cite{global-relay-finra-2026}. More than \$2bn in fines since 2021 for business conversations on unrecorded channels~\cite{american-banker-sec-fines} make any channel described as ``anonymous'' a red flag to compliance teams. And Article 30 of the EU's DORA regulation requires ``unrestricted rights of access, inspection and audit'' over technology providers, plus visibility down the subcontractor chain~\cite{advisera-dora-contractual}. Lethe's firm-held ledger answers the first two, to its credit. The third means an EU firm must register the underlying lab as a subcontractor and seek audit rights that a lab will not grant to a customer it cannot name. A third-party risk review would also see a young vendor holding every prompt in plaintext in place of Microsoft, AWS or Anthropic, for a benefit (identity hiding) that is not on the reviewer's checklist.

Budgets are not the barrier. Two-thirds of Lowenstein's respondents spend over \$1M a year on alternative data, and Harvey's mid-market contracts are estimated at \$300k--\$2M (Vaquill~\cite{vaquill-harvey-pricing}, a low-reliability source). The barrier is that firms spend at that level on capability, and they already get a lab-blind route for free from their cloud. A bank's third-party risk assessment alone can run 90 days (Zycus~\cite{zycus}, a vendor source). Realistic buying cycles are likely 6--18 months at banks and large multi-strategy funds, and 1--3 months at smaller funds where a CTO or COO decides.

Three narrow wedges survive. The first is mid-size funds and smaller firms that cannot get zero retention, relaxed abuse monitoring or Enterprise Frontier Safeguards on their own, since all three require approval or a qualifying agreement. They might buy a sanctioned, pooled zero-retention channel with firm-held logs and access to several models. That is a gateway-as-a-service pitch, with anonymity as a secondary feature. The second is event-driven, activist and M\&A-advisory teams, whose exposure is real but theoretical. Even under zero retention, OpenAI still collects ``account information, metadata, usage statistics, and billing information''~\cite{openai-developer-docs-your}, so a spike in a firm's usage on one sector is visible to the lab when the firm uses the API directly. The third is a lab launching a trading or advisory business, which would turn intent privacy from theory into demand overnight.

\section{Six gaps no competitor closes come down to one: accountable anonymity}

\begin{table}[tbp]
\centering\tablefont
\caption{Six gaps no competitor closes.}\label{tab:gaps}
\begin{tabularx}{\textwidth}{@{}L{1.072}L{1.031}L{1.072}L{0.825}@{}}
\toprule
\colhead{Gap} & \colhead{Closest today} & \colhead{What the labs' terms and safety rules permit} & \colhead{Who has shown willingness to pay} \\
\midrule
Operator cannot link users to frontier-model requests (cryptographically), at commercial quality & OpenAnonymity (academic, free); LLMTor (solo developer); Brave Premium tokens (undocumented) & Permitted as an app serving its own end users (OpenAI \S{}2.2, Anthropic \S{}A.1) & Consumers at the \$15--20 AI price point (Venice, Duck.ai, Brave) \\
Oblivious transport for AI traffic through an independent relay & Confident Security (open-weight only); Apple, Google, Meta (their own models) & No restriction found; relays are a commodity (Cloudflare, Fastly) & Flo Health bought an OHTTP ``Anonymous Mode'' \\
Abuse control and misuse detection across sessions, without knowing identity & LLMTor's permanent ban token; ARC drafts; Tinfoil's in-enclave classifiers (open-weight only) & Labs require detection across sessions but accept monitoring held by the customer (EFS, Private Safety Processing) & The labs, as a condition of access \\
``We don't log'' made checkable at the intermediary & Attested gateways at SayGm, Phala, RedPill & No restriction & Not shown on its own \\
An anonymous tier the provider has formally approved & Duck.ai and Apple contracts (consumer; no Covered Models observed) & Anything resembling resale needs approval (Anthropic \S{}D.4); Mythos limited to approved partners & Not shown \\
Published anonymity-set size and measured residual risk & None & No restriction & Not shown \\
\bottomrule
\end{tabularx}
\tablenote{Sources: as cited in the sections above; Flo Health from Cloudflare~\cite{cloudflare-blog-ohttp-gate}.}
\end{table}

Read across the rows, the six gaps are really one problem: \textbf{anonymous but accountable access}. The labs' binding objection to unlinkable users is not technical. It is safety: Anthropic retains Covered-Model traffic precisely because some misuse shows up only across sessions and accounts. Yet through Enterprise Frontier Safeguards and Private Safety Processing, the labs have already signalled that they will accept monitoring held by the customer rather than by themselves. A pool could be the first intermediary able to give a lab what it needs without telling it who is asking. It would run misuse detection under its own keys, inside an attested environment, without learning who each user is. It would issue anonymous credentials that each allow only a fixed number of requests. And it would give every member a handle that can be suspended but cannot be traced back to a person. There is a telling detail here. Kagi, an intermediary that knows its users through their accounts, lists Claude Fable 5.1, while Duck.ai, which anonymises users under zero retention, stops at Opus 4.8~\cite{kagi-docs-assistant,duckduckgo-help-duck-ai}. That fits, though it does not prove, the idea that Covered Models are reachable only by intermediaries that accept the lab's retention and accountability terms.

\begin{figure}[tbp]
\centering
\figfit{\textwidth}{figs/accountable}
\caption{\textbf{Accountable anonymity.} The labs' binding objection is safety, not technology: some misuse shows up only across sessions. A pool can give a lab that signal without telling it who is asking, by issuing credentials that each allow a fixed number of requests, routing through an independent relay, and running misuse monitoring under its own keys in an attested gateway.}
\label{fig:accountable}
\end{figure}

\section{Verdict: as pitched, Lethe would not be funded}

\subsection{A seed investor or YC would pass on the current pitch}

\begin{callout}[coral]{Verdict}
\textbf{The current strategy will not work as pitched.} It asks investors to believe three things: that the profit centre is a buyer nobody has observed; that the consumer product, which does have proven demand, should run at break-even to subsidise that buyer; and that the moat is crowd size, in a market where DuckDuckGo, Brave, Venice and Apple already have the largest crowds. A YC partner would hear ``dark pool for AI'', ask which fund has asked to hide its name from Anthropic, and find that the written record has no answer. A seed investor would find three diligence failures in the document itself. ``Nobody does this'' is contradicted by OpenAnonymity, LLMTor, Brave, Venice and Duck.ai. ``Most of them restrict frontier models'' is contradicted by JPMorgan, Goldman, Morgan Stanley, Balyasny and Freshfields. And ``a pool with one fund in it hides nobody'' comes with no plan for why the first funds would join a pool that, by its own admission, cannot protect them until it is large. The evidence is not strong enough to say no fund would ever buy, because the written record is silent rather than negative and ten interviews could still surprise. But it is strong enough to say the pitch should not go out in its current form.
\end{callout}

\subsection{The institutional half: the profit centre has no buyer}

The findings against are many, and they reinforce each other. Firms already use frontier models at scale. On Bedrock and Azure the lab already cannot see their data, at no extra cost. The one 2026 shock to frontier data terms was answered with logs held by the customer, while the lab kept knowing who the customer was. No survey, statement or deployment rationale mentions hiding identity. Regulators reward traceability, and DORA demands visibility down the supplier chain. And a startup holding every prompt in plaintext raises vendor risk rather than lowering it.

The labs' safety rules cut deepest. Pooling many firms and an anonymous consumer crowd into one provider account is exactly the pattern that Anthropic's 30-day Covered-Model retention, its approved-partner gating, and OpenAI's suspension and KYC clauses exist to catch. So the best models, the ones institutions most want, are the ones most likely to be withheld. Mixing firms with consumers also creates a new risk for the firms: if a shared key is suspended over one consumer's abuse, every member loses access with it. And the pool cannot hide \emph{what} a firm asks. The UMass study shows text defences fail, and Lermen and colleagues link the same person's pseudonymous profiles at 55\% recall and 90\% precision across 89,000 candidates using content alone~\cite{arxiv-2602-16800-lermen}. A deal team's diligence on one mid-cap company stands out no matter how many consumers sit beside it. Consumers do not fill the anonymity set that matters to a fund.

The findings in favour are real but small. Knowing who looked at what demonstrably carries trading value, as the EDGAR and investor-relations evidence shows, and vendors have abused usage metadata before (Bloomberg, 2013). The labs are entering finance through OpenAI's banker hiring, Claude for Financial Services and Robinhood's agents. On direct API use, firm-level usage metadata survives zero retention. Budgets of \$100k--500k exist. And the labs have shown they will sign with anonymising intermediaries (Duck.ai) and accept monitoring held by the customer (Enterprise Frontier Safeguards). Together these support a future option and a product feature. They do not support a profit centre today.

\subsection{The consumer half: real demand, pitched as a subsidy}

The consumer half is the only part of Lethe with proven demand. Venice is profitable at more than \$70M ARR on an \$18-a-month plan. Duck.ai and Brave have paid tiers, and Kagi, Lumo, Maple and Confer all charge \$10--35 a month. The market pays for privacy-framed AI at the price of AI, not at the price of a privacy tool. The full cryptographic stack has no commercial owner yet, building it is getting cheaper (Cloudflare's gateway, Fastly's relay, ARC), and the labs' terms explicitly permit an app serving its own end users.

The case against is just as concrete. The plan prices this half at break-even. Median first-year revenue per paying subscriber is \$23--32, and only 44\% of annual plans renew~\cite{revenuecat-state-of-subscr}, so at break-even there is nothing left to pay for acquiring users, and the crowd would have to be bought with money the institutional half does not yet earn. It competes with strong incumbents: free Duck.ai, free Apple, Brave's distribution and Venice's 3M users. Its differentiator is invisible to most buyers, because the difference between cryptographic and policy unlinkability does not show on screen. Brave converts about 0.1\% of its 101M monthly users into VPN subscribers~\cite{brave-blog-100m-monthly}, and Mozilla's subscription and advertising revenue together reached only \$66.4M across about 200M users~\cite{mozilla-state-of-mozilla}. OpenAnonymity gives the same primitive away for free and is building the same local-memory ``secretary''. Venice's success has confounders in its ``uncensored'' positioning and its crypto community. And Lethe's claim that it ``cannot'' link users, rather than merely promising not to, does not survive the terms as written. An operator that must suspend specific end users when OpenAI asks (\S{}8.1) needs a handle for each member, and a relay that forwards plaintext to the lab can prove it keeps no logs only if it runs in an attested environment.

\subsection{Four reshaped strategies, ranked}

\begin{table}[tbp]
\centering\tablefont
\caption{Four reshaped strategies, ranked.}\label{tab:strategies}
\begin{tabularx}{\textwidth}{@{}L{0.177}L{1.288}L{1.288}L{1.247}@{}}
\toprule
\colhead{Rank} & \colhead{Strategy} & \colhead{Evidence for} & \colhead{Main risk} \\
\midrule
1 & \textbf{Accountable anonymity as a paid consumer and prosumer product} (detailed below) & Venice, Duck.ai and Brave paid tiers; the full stack has no commercial owner; OpenAnonymity proves it works but has no business; terms permit apps; labs accept customer-held monitoring & Winning users from free incumbents with bigger crowds, on a difference most users cannot see \\
2 & \textbf{License the accountable-anonymity layer} to intermediaries that already pool users (Duck.ai, Brave, Kagi, Venice and similar) and to apps in sensitive categories & Every pooled relay is policy-only; Kagi has not brought Privacy Pass to Assistant; abuse control under anonymity is unsolved; Flo Health paid for OHTTP & Cloudflare, Fastly, ARC and OpenPCC turn it into a commodity; few buyers, and those can build it themselves; infrastructure seed rounds stay small \\
3 & \textbf{Provider-blind, firm-audited frontier access for mid-size firms} (detailed below) & Mid-size firms cannot get zero retention or EFS alone; big firms build this themselves; custody is what banks bought; budgets rising & Crowded (Rogo, Hebbia, Credal, WitnessAI, Bedrock, internal builds); no premium for anonymity, so contract sizes at gateway levels; long security reviews \\
4 & \textbf{Keep the institutional dark pool as an option with a trigger}, not a pitch & Knowing who asked what has value; the labs are moving into finance & The trigger may never come; when it does, lab KYC rules and Covered-Model gating may forbid it \\
\bottomrule
\end{tabularx}
\end{table}

\textbf{Strategy 1 is the strongest, because it keeps the half with proven demand and fixes its economics.} The product offers GPT, Claude and Gemini, plus an attested open-weight tier for content-sensitive work, bought from Phala, NEAR AI or Tinfoil rather than built. Access runs on blind-signature or ARC credentials that each allow a fixed number of requests. Traffic goes through an OHTTP relay run by an independent operator and through a relay running in an attested environment, so that Lethe's own ``we don't log'' can be checked rather than promised. All traffic is padded, the size of the anonymity set is published, and memory stays on the user's device. Lethe itself is the accountable customer: it runs misuse monitoring under its own keys without learning who each user is. That last element is the only credible basis for asking a lab for Duck.ai-style terms. The consumer line stops being the crowd and becomes the business, priced with margin at the \$15--20 that Venice and Duck.ai validated. The first target is the privacy-minded users who already pay Venice, Maple and Proton. For them, the cryptographic difference has to be made visible: a live anonymity-set counter, a public data asset of the kind that drove OpenRouter's growth, and a page stating what the provider, the relay and the token issuer can each still see. The pitch becomes ``Duck.ai with proofs'', not ``dark pool for AI''. It is fundable at seed only with growth numbers. The consumer companies in this space that raised money did so after they had revenue, and the large cheques came from crypto-native funds.

\textbf{Strategy 2 is the natural second act, and a fallback.} If acquiring consumers proves too expensive, the accountable-anonymity layer can be sold to others: token issuance, credentials limited to a fixed number of requests, relay integration, an attested gateway, misuse monitoring and anonymity metrics. It is exactly what Kagi needs to bring Privacy Pass to Assistant, and what relays that rely on policy alone need before they can say ``we cannot''. Its weakness is that Cloudflare co-authors ARC and sells the gateway, Fastly sells the relay, and OpenPCC is open source. The confidential-infrastructure startups in this report raised \$4--5M seeds or died.

\textbf{Strategy 3 is the most conventional path to funding, but it is a different and more crowded company.} It sells what banks have shown they buy. That means a zero-retention channel the provider has approved, a firm-held and immutable ledger that meets SEC Rule 17a-4, monitoring held under the firm's own keys in the style of Enterprise Frontier Safeguards, routing across several models, an attested relay and SOC 2 certification. Masking of firm-level usage metadata becomes a feature for event-driven and advisory teams rather than the headline. The customers are firms too small to negotiate these terms themselves. It needs design partners before any raise, and it should be priced like a gateway or data-loss-prevention product, not as \$250k memberships.

\textbf{Strategy 4 keeps the original idea alive at no cost.} Watch for the triggers, keep the institutional ledger and policy controls in the codebase, and revisit if a lab launches a trading or advisory business. Until then, it is not a pitch.

What to stop doing is just as clear. Stop leading with ``dark pool for AI'' and with \$100k--500k memberships. Stop claiming nobody does this. Stop putting firms and anonymous consumers in the same provider account, and stop treating the consumer app as a subsidy. Do not build standalone open-weight enclave hosting, which Tinfoil, Privatemode, Phala, NEAR AI, Chutes, Maple and Confer already crowd at prices down to \texteuro{}0.65 per million tokens. And do not reach for a token, whose economics in this sector are a warning rather than a template: Nillion has served 1.4M inferences against \$52.7M raised, and Secret's token fell 96\% in a year.

\subsection{Evidence that would reverse this verdict}

Three results would change the institutional judgement, and each can be obtained within the thirty days Lethe's own plan sets aside. The first is three or more funds signing letters of intent that name hiding their identity from the lab, rather than keeping custody of their content, as the reason to buy. The second is one lab stating in writing that it will serve pooled, anonymous institutional traffic, including Covered Models. The third is a third-party risk team at a regulated firm confirming that it could onboard such a pool under DORA and FINRA rules. On the consumer side, the test is whether a landing page that explains cryptographic unlinkability converts privacy-minded users at rates like Venice's. If the interviews turn up only demand for content custody, that is a vote for Strategy 3, not for the pool.

\section{Conclusion}

The research changes what Lethe's problem is. The team saw its risk as cryptography and the cold start. The evidence says the cryptography is the easy part and is becoming a commodity: Apple, Google, Meta, Cloudflare and a Stanford PhD student have all shipped it. The hard part, which nobody has claimed, is accountability: giving a frontier lab the cross-session safety signal it now insists on, without giving it identity. Whoever solves that holds the only contract that matters, because it is the precondition for anonymous access to the best models for anyone, consumer or firm. That is a narrower and more technical moat than ``a bigger crowd'', and unlike a crowd, no incumbent already has it.

The second shift is about who the customer is. In 2026 institutions told the market what they want: custody, audit trails and provider-approved zero retention. They got it from their cloud provider and their lab at no extra cost. Consumers told the market they will pay about \$18 a month for AI that does not profile them, and Venice turned that into a billion-dollar company on policy promises alone. Lethe's original plan sold the cryptographic upgrade to the buyer who did not ask for it and gave it away to the buyer who pays. Reversing that is not a retreat from the institutional vision. It is the only way a pool large enough to protect a fund could ever come to exist.

% ======================================================================  refs
\begin{thebibliography}{999}
\bibitem{arxiv-2609-33903-umass} arXiv 2609.33903 (UMass Amherst study). \url{https://arxiv.org/html/2609.33903v1}
\bibitem{openanonymity-oa-station-r} OpenAnonymity, oa-station-reference (GitHub). \url{https://github.com/OpenAnonymity/oa-station-reference}
\bibitem{openanonymity-oa-chat-priv} OpenAnonymity, oa-chat PRIVACY\_MODEL.md (GitHub). \url{https://github.com/OpenAnonymity/oa-chat/blob/main/docs/PRIVACY_MODEL.md}
\bibitem{openanonymity-oa-verifier} OpenAnonymity, oa-verifier (GitHub). \url{https://github.com/OpenAnonymity/oa-verifier}
\bibitem{openanonymity-oa-chat-gith} OpenAnonymity, oa-chat (GitHub). \url{https://github.com/OpenAnonymity/oa-chat}
\bibitem{ken-liu-on-x} Ken Liu on X. \url{https://x.com/kenziyuliu/status/2027145028737663259}
\bibitem{brave-technologist-podcast} Brave Technologist podcast, episode 120. \url{https://brave.com/podcast/e120/}
\bibitem{the-block-zkapi-launch} The Block (zkAPI launch). \url{https://www.theblock.co/news/defi/2026-10-01-ethereum-foundation-launches-zkapi-417504}
\bibitem{openanonymity-on-github} OpenAnonymity on GitHub. \url{https://github.com/OpenAnonymity}
\bibitem{the-block-venice-funding} The Block (Venice funding). \url{https://theblock.co/post/406934/venice-ai-funding-equity-valuation-dragonfly}
\bibitem{techcrunch-venice-series-a} TechCrunch (Venice Series A). \url{https://techcrunch.com/2026/07/01/venice-ai-becomes-a-unicorn-with-65m-series-a-as-its-privacy-first-ai-platform-takes-off/}
\bibitem{ownyourmind-venice} OwnYourMind (Venice). \url{https://ownyourmind.ai/projects/venice/}
\bibitem{siliconangle-venice} SiliconANGLE (Venice). \url{https://siliconangle.com/2026/07/01/venice-raises-65m-1b-valuation-private-uncensored-ai/}
\bibitem{venice-privacy} Venice, privacy. \url{https://venice.ai/privacy}
\bibitem{venice-docs-privacy} Venice docs, privacy. \url{https://docs.venice.ai/overview/privacy.md}
\bibitem{venice-pricing} Venice, pricing. \url{https://venice.ai/pricing}
\bibitem{duckduckgo-help-duck-ai} DuckDuckGo help, Duck.ai privacy. \url{https://duckduckgo.com/duckduckgo-help-pages/duckai/ai-chat-privacy}
\bibitem{duck-ai-privacy-terms} Duck.ai privacy terms. \url{https://duckduckgo.com/duckai/privacy-terms}
\bibitem{spread-privacy-duckduckgo} Spread Privacy (DuckDuckGo subscription). \url{https://spreadprivacy.com/pro-subscription/}
\bibitem{wikipedia-brave-leo} Wikipedia, Brave Leo. \url{https://en.wikipedia.org/wiki/Brave_Leo}
\bibitem{brave-blog-browser-ai} Brave blog, browser AI in TEEs. \url{https://brave.com/blog/browser-ai-tee/}
\bibitem{kagi-docs-assistant} Kagi docs, Assistant. \url{https://help.kagi.com/kagi/ai/assistant.html}
\bibitem{kagi-docs-privacy-pass} Kagi docs, Privacy Pass. \url{https://help.kagi.com/kagi/privacy/privacy-pass.html}
\bibitem{macrumors-proton-lumo-2} MacRumors (Proton Lumo 2.0). \url{https://macrumors.com/2026/06/30/proton-lumo-version-2-new-features}
\bibitem{apple-chatgpt-extension-pr} Apple, ChatGPT extension privacy. \url{https://www.apple.com/legal/privacy/data/en/chatgpt-extension/}
\bibitem{1privacy-ai} 1privacy.ai. \url{https://1privacy.ai}
\bibitem{nullsink-trust} nullsink, trust. \url{https://nullsink.is/trust/}
\bibitem{anonrouter-ai} anonrouter.ai. \url{https://anonrouter.ai/}
\bibitem{ppq-ai-payperq} ppq.ai (PayPerQ). \url{https://ppq.ai/}
\bibitem{hacker-news-llmtor} Hacker News (LLMTor). \url{https://news.ycombinator.com/item?id=47298201}
\bibitem{tinfoil-docs-attestation-a} Tinfoil docs, attestation architecture. \url{https://docs.tinfoil.sh/verification/attestation-architecture}
\bibitem{launch-hn-tinfoil} Launch HN (Tinfoil). \url{https://news.ycombinator.com/item?id=43996555}
\bibitem{tinfoil-blog-safety-withou} Tinfoil blog, safety without compromising privacy. \url{https://tinfoil.sh/blog/2026-09-14-safety-without-compromising-privacy}
\bibitem{tinfoil-privacy-policy} Tinfoil, privacy policy. \url{https://tinfoil.sh/privacy}
\bibitem{pitchbook-tinfoil} PitchBook (Tinfoil). \url{https://pitchbook.com/profiles/company/770982-94}
\bibitem{tinfoil-blog-tinfoil-conta} Tinfoil blog, Tinfoil Containers. \url{https://tinfoil.sh/blog/2026-03-10-tinfoil-containers}
\bibitem{tinfoil-blog-verifiable-cl} Tinfoil blog, verifiable clean rooms. \url{https://tinfoil.sh/blog/2026-06-09-auditing-frontier-model-without-weights}
\bibitem{confer-blog-private-infere} Confer blog, private inference. \url{https://confer.to/blog/2026/01/private-inference/}
\bibitem{techcrunch-confer} TechCrunch (Confer). \url{https://techcrunch.com/2026/01/18/moxie-marlinspike-has-a-privacy-conscious-alternative-to-chatgpt}
\bibitem{confer-blog-meta-ai} Confer blog, Meta AI. \url{https://confer.to/blog/2026/03/encrypted-meta/}
\bibitem{confer-blog-confidential-w} Confer blog, Confidential Workers. \url{https://confer.to/blog/2026/09/confidential-workers-for-private-ai/}
\bibitem{privatemode-docs-architect} Privatemode docs, architecture. \url{https://docs.privatemode.ai/architecture/overview}
\bibitem{privatemode-pricing} Privatemode, pricing. \url{https://www.privatemode.ai/pricing}
\bibitem{tech-eu-edgeless-systems} Tech.eu (Edgeless Systems seed). \url{https://tech.eu/2023/03/07/edgeless-systems-raises-5million-for-public-cloud-security/}
\bibitem{business-wire-openpcc-laun} Business Wire (OpenPCC launch). \url{https://www.businesswire.com/news/home/20251105013372/en/Confident-Security-Launches-OpenPCC-an-Open-Source-Standard-that-Protects-Data-Shared-with-AI-Models}
\bibitem{finsmes-confident-security} FinSMEs (Confident Security). \url{https://www.finsmes.com/2025/07/confident-security-raises-4-2m-in-funding.html}
\bibitem{techcrunch-confident-secur} TechCrunch (Confident Security). \url{https://techcrunch.com/2025/07/17/confident-security-the-signal-for-ai-comes-out-of-stealth-with-4-2m/}
\bibitem{confident-security} Confident Security. \url{https://confident.security/}
\bibitem{privacy-guides-forum-maple} Privacy Guides forum (Maple). \url{https://discuss.privacyguides.net/t/privacy-using-maple-ai/34503}
\bibitem{maple-blog-maple-proxy} Maple blog, Maple Proxy. \url{https://www.trymaple.ai/blog/introducing-maple-proxy-the-maple-ai-api-that-brings-encrypted-llms-to-your-openai-apps}
\bibitem{human-rights-foundation} Human Rights Foundation. \url{https://hrf.org/latest/hrf-announces-support-for-projects-advancing-ai-for-individual-rights/}
\bibitem{siliconangle-opaque} SiliconANGLE (Opaque). \url{https://siliconangle.com/2026/02/12/opaque-raises-24m-300m-valuation-confidential-ai-platform/}
\bibitem{pr-newswire-opaque-series} PR Newswire (Opaque Series B). \url{https://www.prnewswire.com/news-releases/opaque-raises-24m-series-b-at-300m-valuation-to-advance-confidential-ai-for-the-enterprise-302685635.html}
\bibitem{the-hacker-news-battering} The Hacker News (Battering RAM). \url{https://thehackernews.com/2025/10/50-battering-ram-attack-breaks-intel.html}
\bibitem{tee-fail} TEE.fail. \url{https://tee.fail/}
\bibitem{tinfoil-blog-confidential} Tinfoil blog, confidential computing overhead. \url{https://tinfoil.sh/blog/2026-06-23-confidential-computing-overhead}
\bibitem{near-ai-blog-near} NEAR AI blog, NEAR AI Cloud launch. \url{https://near.ai/blog/introducing-near-ai-cloud-private-chat}
\bibitem{near-ai-blog} NEAR AI blog. \url{https://near.ai/blog}
\bibitem{cryptobriefing-near-ai-git} CryptoBriefing (NEAR AI, GitHub Copilot). \url{https://cryptobriefing.com/near-ai-private-inference-github-copilot/}
\bibitem{near-ai-blog-openrouter} NEAR AI blog, OpenRouter route. \url{https://near.ai/blog/near-ai-cloud-on-openrouter}
\bibitem{cryptobriefing-near-ai-sta} CryptoBriefing (NEAR AI staking). \url{https://cryptobriefing.com/near-ai-staking-confidential-inference/}
\bibitem{coindesk-pha} CoinDesk, PHA. \url{https://www.coindesk.com/price/pha}
\bibitem{phala} Phala. \url{https://phala.com/}
\bibitem{phala-private-ai-inference} Phala, private AI inference. \url{https://phala.com/solutions/private-ai-inference}
\bibitem{phala-cloud-redpill} Phala Cloud, RedPill. \url{https://cloud.phala.network/cases/redpill}
\bibitem{nillion-docs-llms} Nillion docs, LLMs. \url{https://docs.nillion.com/blind-computer/build/llms/overview}
\bibitem{nillion-phase-2-upgrade} Nillion, Phase 2 upgrade. \url{https://nillion.com/news/nillions-phase-2-upgrade-is-live-introducing-a-unified-developer-portal/}
\bibitem{binance-research-nillion} Binance Research, Nillion. \url{https://www.binance.com/en/research/projects/nillion}
\bibitem{secret-network-docs-secret} Secret Network docs, Secret AI architecture. \url{https://docs.scrt.network/secret-network-documentation/secret-ai/architecture}
\bibitem{coingecko-api-scrt} CoinGecko API, SCRT. \url{https://api.coingecko.com/api/v3/coins/secret}
\bibitem{chutes} Chutes. \url{https://chutes.ai/news/private-ai-inference-verifiable-confidential-llm-serving}
\bibitem{iq-wiki-targon} IQ.wiki, Targon. \url{https://iq.wiki/en/wiki/targon}
\bibitem{coingecko-api-pha} CoinGecko API, PHA. \url{https://api.coingecko.com/api/v3/coins/pha}
\bibitem{coingecko-api-nil} CoinGecko API, NIL. \url{https://api.coingecko.com/api/v3/coins/nillion}
\bibitem{tao-media-chutes} TAO Media (Chutes). \url{https://www.tao.media/the-investors-guide-to-chutes-bittensors-inference-layer/}
\bibitem{near-ai-blog-saygm} NEAR AI blog, SayGm. \url{https://near.ai/blog/nearai-cloud-on-saygm-verifiable-confidential-inference}
\bibitem{wisec-2026-arxiv-2605} WiSec 2026, arXiv 2605.24239 (Apple PCC). \url{https://arxiv.org/html/2605.24239v1}
\bibitem{cloudflare-blog-ohttp-gate} Cloudflare blog, OHTTP Gateway. \url{https://blog.cloudflare.com/announcing-cloudflare-ohttp-gateway/}
\bibitem{9to5mac-gemini-and-pcc} 9to5Mac (Gemini and PCC). \url{https://9to5mac.com/2026/01/29/apple-confirms-gemini-powered-siri-will-use-private-cloud-compute/}
\bibitem{apple-security-research-ex} Apple Security Research, expanding PCC. \url{https://security.apple.com/blog/expanding-pcc/}
\bibitem{google-private-ai-compute} Google, Private AI Compute technical brief. \url{https://services.google.com/fh/files/misc/private_ai_compute_technical_brief.pdf}
\bibitem{help-net-security-private} Help Net Security (Private AI Compute memory). \url{https://www.helpnetsecurity.com/2026/09/24/google-private-ai-compute-server-side-memory/}
\bibitem{meta-private-processing-te} Meta, Private Processing technical whitepaper. \url{https://ai.meta.com/static-resource/private-processing-technical-whitepaper}
\bibitem{meta-newsroom-incognito-ch} Meta newsroom, Incognito Chat. \url{https://about.fb.com/news/2026/05/incognito-chat-whatsapp-meta-ai/}
\bibitem{the-register-meta-ai} The Register (Meta AI and ads). \url{https://www.theregister.com/2025/10/01/meta_ai_use_informs_ads/}
\bibitem{openai-help-ads-in} OpenAI help, ads in ChatGPT. \url{https://help.openai.com/en/articles/20001047-ads-in-chatgpt}
\bibitem{tnw-private-safety-process} TNW (Private Safety Processing). \url{https://thenextweb.com/news/openai-zero-data-retention-private-safety-processing}
\bibitem{venturebeat-openai-private} VentureBeat (OpenAI Private Inference). \url{https://venturebeat.com/security/openais-new-private-intelligence-targets-a-growing-enterprise-concern-who-can-see-your-ai-data}
\bibitem{claude-help-covered-models} Claude help, Covered Models data retention. \url{https://support.claude.com/en/articles/15425996-data-retention-practices-for-covered-models}
\bibitem{claude-support-covered-mod} Claude support, Covered Models. \url{https://support.claude.com/en/articles/15425695-covered-models}
\bibitem{anthropic-enterprise-front} Anthropic, Enterprise Frontier Safeguards. \url{https://www.anthropic.com/news/enterprise-frontier-safeguards}
\bibitem{aws-re-post-bedrock} AWS re:Post, Bedrock model data use. \url{https://repost.aws/knowledge-center/amazon-bedrock-model-data-use}
\bibitem{microsoft-learn-azure-open} Microsoft Learn, Azure OpenAI data privacy. \url{https://learn.microsoft.com/en-us/azure/ai-foundry/responsible-ai/openai/data-privacy}
\bibitem{openai-services-agreement} OpenAI Services Agreement. \url{https://cdn.openai.com/osa/openai-services-agreement.pdf}
\bibitem{anthropic-commercial-terms} Anthropic, commercial terms. \url{https://www.anthropic.com/legal/commercial-terms}
\bibitem{anthropic-unsupported-regi} Anthropic, unsupported-region restrictions. \url{https://www.anthropic.com/news/updating-restrictions-of-sales-to-unsupported-regions}
\bibitem{fastly-docs-oblivious-http} Fastly docs, Oblivious HTTP relay. \url{https://docs.fastly.com/products/oblivious-http-relay}
\bibitem{ietf-draft-privacy-pass} IETF draft, Privacy Pass ARC. \url{https://datatracker.ietf.org/doc/html/draft-ietf-privacypass-arc-crypto-01}
\bibitem{yc-oss-yc-companies} yc-oss, YC companies API. \url{https://yc-oss.github.io/api/companies/all.json}
\bibitem{y-combinator-tinfoil} Y Combinator, Tinfoil. \url{https://www.ycombinator.com/companies/tinfoil}
\bibitem{y-combinator-logosguard} Y Combinator, LogosGuard. \url{https://www.ycombinator.com/companies/logosguard}
\bibitem{logosguard-site} LogosGuard site. \url{https://logosguard.com/}
\bibitem{y-combinator-egoist-machin} Y Combinator, Egoist Machines. \url{https://www.ycombinator.com/companies/egoist-machines}
\bibitem{pulse-2-0-unbound} Pulse 2.0 (Unbound). \url{https://pulse2.com/unbound-4-million-raised-for-helping-enterprises-with-ai-tools}
\bibitem{y-combinator-caution} Y Combinator, Caution. \url{https://www.ycombinator.com/companies/caution}
\bibitem{techcrunch-credal} TechCrunch (Credal). \url{https://techcrunch.com/2023/10/26/credal-aims-to-connect-company-data-to-llms-securely/}
\bibitem{y-combinator-blyss} Y Combinator, Blyss. \url{https://www.ycombinator.com/companies/blyss}
\bibitem{y-combinator-litellm} Y Combinator, LiteLLM. \url{https://www.ycombinator.com/companies/litellm}
\bibitem{axios-witnessai} Axios (WitnessAI). \url{https://axios.com/2026/01/13/witnessai-funding-enterprise-ai}
\bibitem{sec-10-q-sentinelone} SEC 10-Q (SentinelOne). \url{https://www.sec.gov/Archives/edgar/data/1583708/000158370825000159/s-20251031.htm}
\bibitem{cyberscoop-cato-aim-securi} CyberScoop (Cato, Aim Security). \url{https://cyberscoop.com/cato-networks-acquires-ai-security-startup-aim-security/}
\bibitem{cb-insights-edgeless-syste} CB Insights, Edgeless Systems. \url{https://www.cbinsights.com/company/edgeless-systems/financials}
\bibitem{techleap-mithril-security} Techleap (Mithril Security). \url{https://finder.techleap.nl/news/feed/h-company-acquires-mithril-security-startup}
\bibitem{contrary-research-openrout} Contrary Research, OpenRouter. \url{https://research.contrary.com/company/openrouter}
\bibitem{bloomberg-stripe-and-openr} Bloomberg (Stripe and OpenRouter). \url{https://www.bloomberg.com/news/articles/2026-08-16/stripe-nears-deal-to-buy-ai-firm-openrouter-for-over-7-billion}
\bibitem{legal-business-harvey} Legal Business (Harvey). \url{https://www.legalbusiness.co.uk/in-house/harvey-hits-15-5bn-valuation-in-largest-fundraising-round-yet/}
\bibitem{sacra-rogo} Sacra, Rogo. \url{https://sacra.com/c/rogo}
\bibitem{let-s-data-science} Let's Data Science (YC Requests for Startups). \url{https://letsdatascience.com/news/yc-requests-ai-native-startups-across-industries-1c422935}
\bibitem{explainx-yc-requests-for} explainx (YC Requests for Startups, Fall 2026). \url{https://explainx.ai/blog/yc-requests-for-startups-fall-2026}
\bibitem{cefpro-jpmorgan-llm-suite} CeFPro (JPMorgan LLM Suite). \url{https://connect.cefpro.com/article/view/inside-jpmorgan-llm-suite-as-ai-agents-spread-across-the-bank}
\bibitem{hargreaves-lansdown-goldma} Hargreaves Lansdown (Goldman Sachs and Anthropic). \url{https://www.hl.co.uk/shares/stock-market-news/company--news/goldman-sachs-working-with-anthropic-to-automate-accounting,-compliance-roles}
\bibitem{openai-morgan-stanley} OpenAI, Morgan Stanley. \url{https://openai.com/index/morgan-stanley/}
\bibitem{cfotech-balyasny} CFOtech (Balyasny). \url{https://cfotech.co.uk/story/balyasny-builds-ai-research-platform-for-hedge-fund-teams}
\bibitem{freshfields-anthropic-part} Freshfields (Anthropic partnership). \url{https://www.freshfields.com/en/our-thinking/news/news-search/2026/04/freshfields-and-anthropic-team-up-to-co-build-ai-legal-workflows-deploying-claude-across-the-firm-globally}
\bibitem{lowenstein-sandler-alterna} Lowenstein Sandler, alternative data report 2025. \url{https://www.lowenstein.com/media/jujd45bp/alt-data-report-2025_final.pdf}
\bibitem{resonanz-capital} Resonanz Capital. \url{https://resonanzcapital.com/insights/ai-use-by-hedge-funds-made-tangible-from-lego-bots-to-alpha-assistants}
\bibitem{aima-hedge-funds-and} AIMA, hedge funds and generative AI. \url{https://www.aima.org/static/a4f9bc40-8c32-42e6-87f52bd89f6e1a82/How-hedge-funds-are-leveraging-Gen-AI-to-get-ahead.pdf}
\bibitem{mayer-brown-m-a} Mayer Brown, M\&A discovery in the AI era. \url{https://www.mayerbrown.com/en/insights/publications/2026/03/ma-discovery-in-the-ai-era-generative-ai-communications-and-outputs-may-become-litigation-ammunition}
\bibitem{management-science} Management Science. \url{https://pubsonline.informs.org/doi/fpi/10.1287/mnsc.2022.4466}
\bibitem{irwin} Irwin. \url{https://getirwin.com/blog/irwin-enhances-website-analytics-solution-irwin-iq-and-expands-access-for-the-investor-relations-community}
\bibitem{cnbc-jpmorgan-and-bloomber} CNBC (JPMorgan and Bloomberg). \url{https://www.cnbc.com/2013/05/16/jpmorgan-demands-bloomberg-staff-logs.html}
\bibitem{cboe-2025-us-equities} Cboe, 2025 US equities year in review. \url{https://www.cboe.com/insights/posts/2025-u-s-equities-year-in-review}
\bibitem{sherwood-news-openai-and} Sherwood News (OpenAI and former bankers). \url{https://sherwood.news/tech/openai-has-an-army-of-ex-investment-bankers-making-financial-models-to-train/}
\bibitem{fortune-robinhood-agents} Fortune (Robinhood agents). \url{https://fortune.com/2026/09/29/robinhood-trading-agents-hood-openai-anthropic/}
\bibitem{global-relay-finra-2026} Global Relay (FINRA 2026 oversight report). \url{https://www.globalrelay.com/resources/thought-leadership/finra-2026-oversight-report-flags-genai-recordkeeping-and-cybersecurity-risks/}
\bibitem{american-banker-sec-fines} American Banker (SEC fines). \url{https://www.americanbanker.com/news/judge-orders-sec-to-release-data-behind-2b-in-whatsapp-fines}
\bibitem{advisera-dora-contractual} Advisera (DORA contractual provisions). \url{https://advisera.com/dora-regulation/key-contractual-provisions/}
\bibitem{vaquill-harvey-pricing} Vaquill (Harvey pricing). \url{https://www.vaquill.ai/blog/harvey-ai-price-increase-per-seat-2026}
\bibitem{zycus} Zycus. \url{https://www.zycus.com/?p=205568}
\bibitem{openai-developer-docs-your} OpenAI developer docs, your data. \url{https://developers.openai.com/api/docs/guides/your-data}
\bibitem{arxiv-2602-16800-lermen} arXiv 2602.16800 (Lermen et al.). \url{https://arxiv.org/abs/2602.16800}
\bibitem{revenuecat-state-of-subscr} RevenueCat, State of Subscription Apps 2025. \url{https://www.revenuecat.com/state-of-subscription-apps-2025/}
\bibitem{brave-blog-100m-monthly} Brave blog, 100M monthly users. \url{https://brave.com/blog/100m-mau/}
\bibitem{mozilla-state-of-mozilla} Mozilla, State of Mozilla. \url{https://stateof.mozilla.org/}
\end{thebibliography}

\end{document}
